\documentclass[../../main.tex]{subfiles}
% \documentclass[a4paper,10pt]{ltjsarticle}
\begin{document}
\rhead{京都大学後期$1995$年$3$番}
\setcounter{equation}{0}
\begin{oframed}
  \documentclass[../../main.tex]{subfiles}
\begin{document}

$a$，$b$，$c$は実数で$a\ge 0$，$b\ge 0$とする．
\begin{align*}
  p(x) & =ax^2+bx+c \\
  q(x) & =cx^2+bx+a
\end{align*}
とおく．$-1\le x\le 1$を満たす全ての$x$に対して$|p(x)|\le 1$が成り立つ時，$-1\le x\le 1$を満たす全ての$x$に対して$|q(x)|\le 2$が成り立つことを示せ．

\end{document}

\end{oframed}
\setlength{\columnseprule}{0.4pt}
\begin{multicols}{2}
  {\bf [解答]}

  $c=0$か否かで場合わけして考える．

  まず，$c=0$の時，
  $q(x)$ は高々1次関数で，その最大値・最小値は端点で取るので$|x| \leqq 1$ の時，
  \begin{align}
    & \min\{q(1), q(-1)\} \leqq q(x) \leqq \max\{q(-1), q(1)\} \\
    \iff
    & a-b \leqq q(x) \leqq a+b \label{eq:1}
  \end{align}
  である．ただし，$a,b \ge 0$であることを用いた．

  一方で，$|p(\pm 1)| \leqq 1$ だから
  \begin{align}
    -1 \leqq a \pm b \leqq 1  \label{eq:2}
  \end{align}
  である．

  \cref{eq:1,eq:2}から
  \begin{align}
    -1 \le q(x) \le 1
  \end{align}
  となって$|q(x)| \le 2$は成立する．

  次に $c \neq 0$ とする．
  $q(x)$ は2次関数で，その最大値・最小値は軸ないし端点でとる．
  軸は$x=-\dfrac{b}{2c}$だから，これらの点での$q(x)$の値は
  \begin{align}
    \begin{dcases}
      q(1) = a+b+c  \\
      q(-1) = a-b+c \\
      q\left(-\dfrac{b}{2c}\right) = a-\dfrac{b^2}{4c}
    \end{dcases}
  \end{align}
  で，これらの絶対値が2以下ならば題意の条件を満たす．

  ここで，$|p(x)| \le 1$より
  \begin{align}
    |p(\pm 1)| = |a \pm b + c| \leqq 1
  \end{align}
  であるから，
  \begin{align}
    |q(\pm 1)| = |p(\pm 1)| \le 1 \label{eq:3}
  \end{align}
  となる．

  次に軸での値について，軸が最大または最小となりうるのは軸が$-1\le x\le 1$の範囲にあるときで
  \begin{align}
    & -1 \le \dfrac{-b}{2c} \le 1                        \\
    \iff
    & -1 \leqq \dfrac{b}{2c} \leqq 1 \label{eq:4}
  \end{align}
  の場合を考えれば良い．この時三角不等式より
  \begin{align}
    \left|q\left(-\dfrac{b}{2c}\right)\right|
    & = \left|a-\dfrac{b^2}{4c}\right|                                                 \\
    & \leqq |a| + |\dfrac{b^2}{4c}|                                                    \\
    & \leqq |a| + \left|\dfrac{b}{2}\right|\left|\dfrac{b}{2c}\right|                   \\
    & \leqq |a| + \left|\dfrac{b}{2}\right| \quad (\because \text{\cref{eq:4}}) \\
    & \leqq |a+b| \quad (\because b \ge 0) \label{eq:5}
  \end{align}
  となる．

  一方で，$p(x)$の条件と三角不等式より
  \begin{align}
    & |p(1)-p(0)| \le |p(1)| + |p(0)| \le 2 \\
    \therefore
    & |a+b| \le 2 \label{eq:6}
  \end{align}
  だから，\cref{eq:5,eq:6}より
  \begin{align}
    \left|q\left(-\dfrac{b}{2c}\right)\right| \le 2 \label{eq:7}
  \end{align}
  である．

  以上\cref{eq:3,eq:7}より，全ての$x, |x|\leqq 1$に対して
  $|g(x)| \leqq 2$ である．
  よって $c \neq 0$ の場合も示された．

  以上で全ての場合が尽くされ，題意は示された．$\cdots$(答)

  \vspace{10pt}
  {\bf [解説]}
  解答では丁寧に場合分けして解いたが，もっとエレガントに条件を駆使して解けるので紹介する．
  ポイントはうまいこと$q(x)$を$p(x)$を利用して表すことである．
  三角不等式から
  \begin{align}
    |q(x)|
    & = |cx^2+bx+a|               \\
    & = |c(x^2-1) + a+bx+c|       \\
    & \leqq |c||x^2-1| + |bx+a+c|
  \end{align}
  であるが，$-1\le x \le 1$の範囲で
  \begin{align}
    & |c||x^2-1| \le |c| = |p(0)| \le 1           \\
    & |bx+a+c| \le \max(|\pm b+a+c|) = |p(\pm 1)| \le 1
  \end{align}
  だから，
  \begin{align}
    |q(x)|
    & \leqq |c||x^2-1| + |bx+a+c| \\
    & \leqq 1 + 1                 \\
    & = 2
  \end{align}
  となり題意は示された．

  \newpage
\end{multicols}
\end{document}
